\scanpage{173}
\mspage{491}
% position: upper paragraph
contact avec des choses décidément juteuses - chacun de ces groupes était ri-\linebreak che de poids et de mystère, et chacun des couples yin-yang qui étaient censés le constituer (mais qui plutôt, tous ensemble, le \mbox{d é s i g n e n t}, sans aucunement l'épuiser) - chacun de ces couples a quelque chose de délicat et d'important à me dire sur la nature de ce monde dans lequel je vis, et souvent sur ma propre nature. J'ai retrouvé avec une force nouvelle ce sentiment qui était présent déjà il y a cinq ans: que le jeu délicat du yin et du yang, du ``féminin'' et du ``mâle'' en toutes choses, est un fil conducteur incomparable vers une compréhension du monde et de soi. Il nous conduit droit vers les ques-\linebreak tions essentielles. Souvent aussi, le ``yoga'' même du yin et du yang, le seul fait\inserted{\unsure{,}} j'entends\inserted{\unsure{,}} de faire attention à l'aspect des choses et évènements qui s'ex-\linebreak prime en termes d'équilibre et de déséquilibre yin-yang, fournit une première clef pour une meilleure compréhension de ces questions, et vers une réponse.

% position: middle paragraph
Je m'excuse si pour certains lecteurs je dois donner l'impression, depuis une page ou deux, de parler du sexe des anges, alors qu'ils ne verraient pas trop même quels sont ces fameux ``couples'' yin- yang dont je parle, et encore moins ces ``groupes'' en lesquels certains se réunissent, lesquels groupes fina-\linebreak lement seraient censés s'assembler en un ``diagramme'' (c'est quand même utile les maths!). Je devrais donner ici au moins un de ces groupes - et j'ai envie de prendre celui par lequel spontanément j'avais commencé hier, celui aussi qui a fini par apparaître au cours de la réflexion comme le groupe ``primitif'' (*), dont semblent sortir progressivement tous les autres, par des sortes de ``filia-\linebreak tions'' successives (se poursuivant sur mon fameux diagramme sur huit ``généra-\linebreak tions''...). Voici donc la liste des ``couples'' que j'ai relevés, constituant ce groupe primitif (qu'on pourrait nommer par le premier de ces couples, savoir ``le groupe \mbox{a c t i o n - i n a c t i o n}'').
% position: lower enumeration
\diagram{$\begin{array}{l}
\text{Action-inaction}\\
\text{activité-passivité}\\
\text{veille-sommeil}\\
\text{sujet-objet}\\
\text{engendrer-concevoir (*)}\\
\text{exécution-conception (*)}
\end{array}$}
% position: footnote rule
\par\noindent\rule{0.25\linewidth}{0.4pt}
% position: footnote
(*) (6 novembre) En fait , il y a un groupe plus primitif encore, qu'on peut appeler le groupe ``\mbox{p è r e - m è r e}''. Voir, au sujet de cet ``oubli'', la note ``Notre Mère la Mort - ou l'Acte et le tabou'' (n° 113). Les couples ``engendrer-concevoir'' et``exécution-conception'', que j'ai fait figurer ci-dessous dans le groupe (soi-disant ``primitif'') \unsure{\inserted{ac}}tion-inaction, s'insèrent visiblement de façon plus naturelle dans le ``groupe mère'' formé autour du couple ``père-mère''.
% position: foot pencil
\margin{172}
